\begin{summarybox}
\textbf{\textcolor{CSummary}{一句话核心}}
切点弦是\textbf{两条切线的两个切点的连线}；其方程来自“\textbf{两个切点处的切线都经过同一点 $P$}”，不是生搬硬套换字母。

\medskip
\textbf{\textcolor{CSummary}{先判断，再写公式}}
\[
\begin{aligned}
 y^2=2px,\quad y_0^2-2px_0>0
 &\quad\Longrightarrow\quad \boxed{y_0y=p(x+x_0)},\\[0.4em]
 x^2=2py,\quad x_0^2-2py_0>0
 &\quad\Longrightarrow\quad \boxed{x_0x=p(y+y_0)}.
\end{aligned}
\]
这里均假设 $p>0$。判断值 $=0$ 时只有一个切点，$<0$ 时没有实切点。

\medskip
\textbf{\textcolor{CSummary}{四种证明：学习与考场的不同用途}}
\begin{enumerate}[leftmargin=*,itemsep=0.3em]
\item \textbf{门槛最低（学生自推）：} 两切点处切线都经过 $P$，\textbf{相减求斜率、相加找中点关系、点斜式还原}。
\item \textbf{速度最快：} 相减求斜率后，直接用切点 $A$ 写点斜式，不必引入中点。
\item \textbf{步骤最少：} 直接证明两个切点都满足 $y_0y=p(x+x_0)$，不需要求斜率。
\item \textbf{常规通用：} 设 $x=x_0+m(y-y_0)$，联立抛物线，用判别式为零求切点并检验所在直线。
\end{enumerate}

\textbf{相减与相加的几何意义（$y_0\ne0$）}
\[
 \underbrace{y_0(y_1-y_2)=p(x_1-x_2)}_{\text{相减：确定弦方向}}
 \quad\Rightarrow\quad k_{AB}=\frac p{y_0},
\]
\[
 \underbrace{y_0y_m=p(x_0+x_m)}_{\text{相加：建立中点关系}}
 \quad\Rightarrow\quad
 y-y_m=\frac p{y_0}(x-x_m)
 \quad\Rightarrow\quad y_0y=p(x+x_0).
\]
当 $y_0=0$ 时，弦为 $x=-x_0$，不能除以 $y_0$。\textbf{相加得到的是中点关系，不是已求出中点坐标。}

\medskip
\textbf{\textcolor{CSummary}{高价值推论}}
若 $y^2=2px$ 的切点弦中点为 $M$，则
\[\boxed{M\left(\frac{y_0^2}{p}-x_0,\ y_0\right)}.\]
切点纵坐标 $y_1,y_2$ 是 $t^2-2y_0t+2px_0=0$ 的两个根，故 $y_1+y_2=2y_0$：\textbf{弦中点与 $P$ 同高，但 $P$ 通常不在弦上}。

\medskip
\textbf{\textcolor{CSummary}{考场三步自查}}
\textbf{辨平方字母 $\to$ 算存在条件 $\to$ 写弦方程（需要切点再联立）}。
\end{summarybox}
